% Generated draft: requires mathematical and source review.
\documentclass[11pt]{article}
\usepackage[margin=25mm]{geometry}
\usepackage{fontspec}
\setmainfont{Latin Modern Roman}
\usepackage{amsmath,amssymb,mathtools,unicode-math}
\setmathfont{Latin Modern Math}
\usepackage{enumitem,needspace,xurl,hyperref}
\setlength{\parindent}{0pt}
\setlength{\parskip}{5pt}
\setlength{\emergencystretch}{3em}
\newcommand{\N}{\mathbb{N}}
\newcommand{\Z}{\mathbb{Z}}
\newcommand{\Q}{\mathbb{Q}}
\newcommand{\R}{\mathbb{R}}
\newcommand{\C}{\mathbb{C}}
\newcommand{\F}{\mathbb{F}}
\newcommand{\A}{\mathbb{A}}
\newcommand{\PP}{\mathbb{P}}
\newcommand{\E}{\mathbb{E}}
\newcommand{\eps}{\varepsilon}
\newcommand{\dd}{\,\mathrm{d}}
\begin{document}
\section*{UnsolvedMath ID: 20001992}
\textbf{Source problem code:} catalogue entry
\subsection{Definitions and mathematical statement}
Eugene Tevelev and Sean Keel ask for the precise relationship between the tropical Grassmannian and Kapranov's Chow quotient \(\operatorname{Gr}(r,n)//(\C^*)^n\), motivated by the cases \((2,n)\) and \((3,6)\) studied by Speyer--Sturmfels. The double slash here denotes a Chow quotient, not an unspecified GIT quotient at a chosen linearization.

For \(1\leq r<n\), let \(N=\binom nr\) and let \(\operatorname{Gr}(r,n)^\circ\) be the all-Plücker-nonzero open subset. Column scaling has effective torus
\[
 T_{\mathrm{eff}}=(\C^*)^n/\C^*_{\mathrm{diag}},\qquad
 X=\operatorname{Gr}(r,n)^\circ/T_{\mathrm{eff}},
\]
embedded in the quotient of the projective Plücker torus. In affine Plücker coordinates the real lineality to remove is
\[
 L=\left\{\left(\sum_{i\in I}a_i\right)_{|I|=r}:a\in\R^n\right\}.
\]
Thus use \(\operatorname{Trop}(\operatorname{Gr}(r,n))/L\), or the equivalent projective quotient by the \((n-1)\)-dimensional effective column lineality. Removing only the overall scalar line leaves unwanted directions. Integral toric constructions also require the corresponding saturated quotient lattice. The resulting open quotient has dimension \(r(n-r)-n+1\).

Choose a fan \(\Sigma\) supported on this tropical quotient and a compatible toric compactification \(\overline X\subset Y_\Sigma\), with flat surjective multiplication map
\[
 S\times\overline X\longrightarrow Y_\Sigma,
\]
where \(S\) is the ambient quotient torus. For a \(k\)-dimensional cone \(\sigma\in\Sigma\), the aid's orbit-codimension test is
\[
 \overline X\cap O(\sigma)\ne\varnothing,
 \qquad\overline X\cap O(\sigma)
 \text{ has pure codimension }k\text{ in }\overline X.
\]
In particular its dimension is \(\dim X-k\). This is not by itself a bijection between cones and irreducible boundary components: intersections may have several components, and choices of fan refinement matter. With the compatible Chow/toric models described by the aid, the identity on the open quotient extends birationally to the Chow compactification; such an extension is not claimed for every unrelated fan.

For rank two, \(n\geq4\), the supplied relationship is the isomorphism
\[
 \overline X\simeq\overline M_{0,n}
 \simeq\operatorname{Gr}(2,n)//T_{\mathrm{eff}}.
\]
After the lineality quotient, cones are labelled by trees with \(n\) labelled leaves and nonnegative internal-edge lengths. An internal edge determines a split \(A\mid A^c\) with \(|A|,|A^c|\geq2\), giving boundary divisor \(D_A=D_{A^c}\). A tree with \(k\) internal edges gives a \(k\)-cone and a boundary stratum of codimension \(k\). Compatible splits describe the intersections. Face/closure order is reversed: if \(\tau\) is a face of \(\sigma\), the stratum for \(\sigma\) lies in the closure of the stratum for \(\tau\).

For \((r,n)=(3,6)\), the aid instead reports a tropical compactification giving a crepant resolution of a Chow quotient with forty isolated singularities. Its dimension is four. Thus its tropical fan describes boundary incidence on the resolved model, rather than a one-to-one Chow-boundary stratification. Crepant means equality of the canonical divisor with the pullback from the quotient in the appropriate Cartier framework. The exact fan/refinement, normalization of the Chow quotient, canonical-divisor hypotheses and reference are not supplied, so this report is restricted to that intended model and is not asserted for every tropical fan on the same support. No complete relationship in arbitrary \((r,n)\) is established by these two cases.
\subsection{Short English statement}
The relationship requires removing the full Plücker lineality and specifying a compactifying fan. In rank two the tree fan matches the boundary of \(\overline M_{0,n}\). For \((3,6)\), the aid reports that the tropical model resolves the Chow quotient, so its cones record the resolved boundary rather than the quotient boundary directly.
\subsection{Sources}
\begin{itemize}
\item Permanent UnsolvedMath ID \(20001992\), exact public code catalogue entry; local registry \texttt{aim-geometry-notes.json}, \texttt{notes[329]}, canonical tag section.
\item AIM workshop ``Amoebas and tropical geometry,'' item A.10, question attributed to Eugene Tevelev and Sean Keel, contributed by Sean Keel: \url{https://aimath.org/WWN/amoebas/amoebas.pdf}. The supplied verification report identifies a PDF version dated January 14, 2004 and says it was checked directly during that research run; no PDF was accessed for this draft.
\item Original attributions: Speyer--Sturmfels for tropical Grassmannians and Kapranov for the Chow quotient. The aid additionally invokes Tevelev's tropical compactification. Full titles, theorem citations and the reference for the forty-singularity crepant model are absent.
\item Research attempt 1, ``From tropical Grassmannians to Chow quotients,'' labelled partial\_result, partially\_solved, proved and low novelty confidence. No original literature entry or proposer/year metadata is supplied beyond the in-text attribution.
\end{itemize}
\end{document}
